\section{Critical label-provenance confound in the UniProt scale-up}
\label{sec:labelconfound}
The positive UniProt REST query was intended to request reviewed KW-0929
entries, but the actual recorded URL lacks the \texttt{reviewed:true}
constraint. The raw positive FASTA contains 27{,}392 entries, of which
24{,}295 have the UniProt \texttt{tr|} prefix (unreviewed) and only 3{,}098
have \texttt{sp|} (reviewed). After canonical-residue filtering and
sequence de-duplication, 24{,}628 positives remained: 21{,}742 unreviewed
and 2{,}886 reviewed. The negative pool was queried with
\texttt{reviewed:true}, so all 24{,}628 selected negatives are reviewed.
This is a major source-status confound, not an innocuous metadata issue.

\begin{table}[h]\centering
\caption{Review status by class and split, counted from the committed CSV
headers. \texttt{tr|} indicates UniProtKB/TrEMBL; \texttt{sp|} indicates
UniProtKB/Swiss-Prot.}
\label{tab:reviewconfound}
\begin{tabular}{lrrrr}
\hline
Split & AMP $\mathrm{tr}$ & AMP $\mathrm{sp}$ & Non-AMP $\mathrm{tr}$ & Non-AMP $\mathrm{sp}$ \\
\hline
Train & 19{,}173 & 2{,}376 & 0 & 19{,}227 \\
Validation & 1{,}347 & 271 & 0 & 2{,}476 \\
Test & 1{,}222 & 239 & 0 & 2{,}925 \\
\hline
Total & 21{,}742 & 2{,}886 & 0 & 24{,}628 \\
\hline
\end{tabular}
\end{table}

The models do not receive review status as an explicit input, but annotation
pipelines, organism composition, length distributions within buckets, and
family composition can carry review-status proxies. The 0.921 AUROC for
PepCNN on the cluster-aware test set is therefore a measurement on this
\emph{confounded} dataset, not a reliable estimate of AMP recognition on a
review-status-matched population. Ten-mer clustering blocks exact motif
sharing across splits; it does not remove a train/test review-status
association. Likewise, APD6 validation with external positives and the
same reviewed test negatives cannot fully repair a training confound.

The proper repair is to fetch reviewed-only positives or select negatives
matched on review status and organism group, rebuild all splits, retrain,
and re-evaluate before claiming a clean scale-up benchmark. The small pilot
FASTA files in the handoff have synthetic identifiers (u0/n0), so their
per-record source and review status cannot be independently verified from
the shipped fixture; the pilot is a software sanity check, not a provenance-
controlled published benchmark. We preserve the
original experiment as an auditable negative result rather than delete it.
The original source query and exact split counts are in
\texttt{data/raw/large\_provenance.json} and
\texttt{results/large\_label\_provenance.json}. The Veltri 2018 and
Feb2020 experiments use separate public benchmark corpora and are not
contaminated by this UniProt query error, though they have their own
random-split homology caveat.
